\section{System overview}
\label{sec:system}

\subsection{The card}
The DU-PCB-001-003 carries a Kintex-7 \code{xc7k480tiffg1156-2L} (299\,k LUTs, 597\,k flip-flops, 955 BRAM36 tiles,
1920 DSP48E1 slices) on an x8 PCIe edge connector, with two independent DDR3 channels. Each channel has eight
2\,Gbit \x8 chips (64-bit data, 2\,GB) plus a ninth chip on an ECC lane that is not used. Two 200\,MHz oscillators,
a 50\,MHz oscillator and the 100\,MHz PCIe reference clock are available. The design is loaded over JTAG
(writing the configuration flash is not supported on this card). In the development PC the slot trains the
link to Gen2 x4.

\subsection{The SoC}
The base LiteX SoC for the card (\file{boards/dupcb/target.py}) has a VexRiscv CPU, both DDR3 controllers and the
PCIe endpoint. The accelerator uses a variant without the CPU, \code{GemmSoC} (\file{hw/gemm/gemm\_soc.py}),
shown in \cref{fig:system}:
\begin{itemize}
  \item \textbf{No CPU.} Both channels (\(2\times2\)\,GB) belong to the host DMA and the kernel. The host does
        everything a CPU would do: DDR3 calibration, kernel control, reading counters.
  \item \textbf{BAR0 is the CSR bus.} The host reads and writes the SoC's 32-bit control and status registers
        (CSRs) through a 128\,KB PCIe memory window. Their addresses come from the build's generated \code{csr.h}.
  \item \textbf{One DMA channel}, used by a bridge that moves data between the PCIe DMA streams and either DDR3
        channel (\cref{sec:pcie}).
  \item \textbf{The HLS kernel} with three AXI4 masters (A, B, C). Each is converted to a LiteDRAM native port on the
        channel the build's \code{memory\_map} assigns to that matrix (\cref{sec:kernel}).
  \item \textbf{Identification.} A \code{build\_id} register holds the first 32 bits of the build's configuration
        hash, so the host can refuse to run against the wrong bitstream.
\end{itemize}
Channel 0 keeps LiteX's standard \code{add\_sdram} integration (its Wishbone port stays idle without a CPU), so the
DDR register names, and with them the calibration script, are the same as in the base SoC. Without the kernel the
SoC uses 25.5\,k LUTs, 19.5\,k flip-flops and 73.5 BRAM tiles.

\begin{figure}[t]
  \centering
  \resizebox{\linewidth}{!}{% System overview: host PC, PCIe link, the CPU-less GemmSoC and both DDR3 channels.
\begin{tikzpicture}[x=1cm,y=1cm, every node/.append style={font=\sffamily\footnotesize}]
  \tikzset{w/.style={minimum width=#1}}
  % ---------------- Host PC ----------------
  \node[hst, w=31mm] (app)   at (0, 0)    {\code{gemm\_runner} (C++17)\\[-1pt]\scriptsize Python: \code{gemm\_run}, \code{gemm\_dse}};
  \node[hst, w=31mm] (lib)   at (0,-1.3)  {host library\\[-1pt]\scriptsize \code{Card} · \code{DdrDma} · \code{GemmKernel}};
  \node[hst, w=31mm] (llp)   at (0,-2.5)  {\code{liblitepcie}};
  \node[hst, w=31mm] (drv)   at (0,-3.9)  {\code{litepcie.ko}\\[-1pt]\scriptsize ioctl · mmap · MSI};
  \node[hst, w=31mm] (ring)  at (0,-5.2)  {DMA buffers\\[-1pt]\scriptsize H2C + C2H, 256\x8\,KB each};
  \node[hst, w=31mm] (iommu) at (0,-6.5)  {IOMMU · root port};
  \draw[link] (app) -- (lib);
  \draw[link] (lib) -- (llp);
  \draw[biln] (llp) -- node[lblw, left=0.5mm] {\code{/dev/litepcie0}} (drv);
  \draw[biln] (drv) -- (ring);
  \draw[biln] (ring) -- (iommu);
  \draw[densely dotted, black!50] (-1.75,-3.2) -- (1.75,-3.2);
  \node[lbl, anchor=south east] at (1.75,-3.2) {user space};
  \node[lbl, anchor=north east] at (1.75,-3.2) {kernel};
  \begin{scope}[on background layer]
    \node[region=black, fit=(app)(iommu), inner sep=6pt] (host) {};
  \end{scope}
  \node[regionl=black, anchor=south west] at (host.north west) {Host PC (Linux, x86-64)};

  % ---------------- FPGA ----------------
  \def\cA{5.6} \def\cB{8.85} \def\cC{12.1} \def\cD{15.35}
  \node[hard, w=27mm] (hb)  at (\cA, 0)    {PCIe hard block\\[-1pt]\scriptsize \code{PCIE\_2\_1} + 8 GTX};
  \node[pci,  w=27mm] (phy) at (\cB, 0)    {PCIe PHY\\[-1pt]\scriptsize \code{S7PCIEPHYRXBE}, CDC};
  \node[sys,  w=27mm] (ep)  at (\cC, 0)    {LitePCIe endpoint\\[-1pt]\scriptsize TLP layer, crossbar};
  \node[sys,  w=27mm] (dma) at (\cD, 0)    {DMA reader + writer\\[-1pt]\scriptsize descriptor tables, MSI};
  \node[sys,  w=27mm] (csr) at (\cC,-1.6)  {BAR0 \(\rightarrow\) CSR bus\\[-1pt]\scriptsize 128\,KB window};
  \node[sys,  w=27mm] (br)  at (\cD,-1.6)  {DDR DMA bridge\\[-1pt]\scriptsize 128\(\leftrightarrow\)512\,b, channel sel.};
  \node[acc,  w=38mm, minimum height=15mm] (k) at (\cA+0.55,-3.55)
       {HLS GEMM kernel\\[-1pt]\scriptsize chain of \(P_N\) PEs, \(P_N P_M\) MAC/cycle\\[-1pt]\scriptsize \code{m\_axi} A, B, C (512\,b)};
  \node[acc,  w=20mm] (cvt) at (\cB+0.35,-3.55) {AXI\(\rightarrow\)native\\[-1pt]\scriptsize + perf taps};
  \node[sys,  w=27mm] (xb0) at (\cC,-3.55)  {crossbar CH0\\[-1pt]\scriptsize CDC, 512\,b ports};
  \node[sys,  w=27mm] (xb1) at (\cD,-3.55)  {crossbar CH1\\[-1pt]\scriptsize CDC, 512\,b ports};
  \node[sys,  w=27mm] (mc0) at (\cC,-5.25)  {LiteDRAM controller\\[-1pt]\scriptsize 8 bank machines};
  \node[sys,  w=27mm] (mc1) at (\cD,-5.25)  {LiteDRAM controller\\[-1pt]\scriptsize 8 bank machines};
  \node[ddr,  w=27mm] (ph0) at (\cC,-6.5)   {\code{WL7DDRPHY} CH0\\[-1pt]\scriptsize 1:4, CK phase by MMCM};
  \node[ddr,  w=27mm] (ph1) at (\cD,-6.5)   {\code{WL7DDRPHY} CH1\\[-1pt]\scriptsize 1:4, CK phase by MMCM};
  \node[blk, fill=black!4, w=38mm, align=left, font=\sffamily\scriptsize] (misc) at (\cA+0.55,-6.0)
       {\textbf{also on the CSR bus}\\ kernel start/args/status, perf counters,\\ \code{build\_id}, DDR calibration + BIST,\\
        bridge control, freq.\ meters, LEDs};

  \draw[wbi] (hb) -- (phy);
  \draw[wbi] (phy) -- (ep);
  \draw[wbi] (ep) -- (dma);
  \draw[biln] (ep) -- node[lblw, left] {MMIO} (csr);
  \draw[wbi] (dma) -- node[lblw, right] {128\,b stream} (br);
  \draw[wbi] ([xshift=5mm]br.south) -- ([xshift=5mm]xb1.north);
  \draw[wbi] ([xshift=-5mm]br.south) -- ++(0,-0.55) -| ([xshift=5mm]xb0.north);
  \draw[wbi] (k) -- (cvt);
  \draw[wbi] (cvt) -- node[lbl, above] {A} (xb0);
  \draw[wbi] (cvt.south) -- ++(0,-0.62) -| node[lbl, pos=0.12, below] {B, C} ([xshift=-6mm]xb1.south);
  \draw[wbi] (xb0) -- (mc0);  \draw[wbi] ([xshift=3mm]xb1.south) -- ([xshift=3mm]mc1.north);
  \draw[wbi] (mc0) -- (ph0);  \draw[wbi] (mc1) -- (ph1);
  \draw[ctl] (csr.west) -| node[lblw, pos=0.25, above] {start, arguments, status, counters} (k.north);

  \begin{scope}[on background layer]
    \node[region=cGrey, fit=(hb)(dma)(k)(misc)(ph0)(ph1), inner sep=7pt] (fpga) {};
  \end{scope}
  \node[regionl=cGrey, anchor=south west] at (fpga.north west) {Kintex-7 480T (\code{xc7k480tiffg1156-2L}): CPU-less GemmSoC, LiteX 2026.08};

  % ---------------- DDR3 and the link ----------------
  \node[ddr, w=27mm] (d0) at (\cC,-8.25) {DDR3 CH0 · 2\,GB\\[-1pt]\scriptsize 8 chips \x8, 64\,b data};
  \node[ddr, w=27mm] (d1) at (\cD,-8.25) {DDR3 CH1 · 2\,GB\\[-1pt]\scriptsize 8 chips \x8, 64\,b data};
  \draw[wbi] (ph0) -- node[lbl, right] {1000\,MT/s} (d0);
  \draw[wbi] (ph1) -- node[lbl, right] {1000\,MT/s} (d1);
  \draw[wbi, draw=cGreen!55!black] (iommu.east) -- (2.55,-6.5) -- node[lbl, sloped, above, pos=0.5]
       {PCIe Gen2 · 5\,GT/s per lane · x4 in this slot (x8 card)} (2.55,0) -- (hb.west);

  % legend
  \node[sys, minimum height=4mm, font=\sffamily\scriptsize] (l1) at (0.3,-8.25) {sys 125\,MHz};
  \node[acc, minimum height=4mm, font=\sffamily\scriptsize, right=1.5mm of l1] (l2) {accel 150--250\,MHz};
  \node[pci, minimum height=4mm, font=\sffamily\scriptsize, right=1.5mm of l2] (l3) {PCIe clocks};
  \node[ddr, minimum height=4mm, font=\sffamily\scriptsize, right=1.5mm of l3] (l4) {DDR PHY/DRAM};
  \node[hard, minimum height=4mm, font=\sffamily\scriptsize, right=1.5mm of l4] (l5) {hard IP};
\end{tikzpicture}
}
  \caption{System overview. The host talks to the card over PCIe: memory-mapped CSR accesses for control and two
    DMA directions for data. On the FPGA, the DMA bridge and the kernel share two LiteDRAM crossbars, one per DDR3
    channel. Colours mark the clock domains.}
  \label{fig:system}
\end{figure}

\subsection{Clocks}
\Cref{fig:clocks} shows where the clocks come from and which logic runs on each. The \textbf{sys} clock
(125\,MHz, from oscillator \#0) runs the CSR bus, the LitePCIe endpoint and DMA, the bridge and both LiteDRAM
controllers. Each DDR3 PHY has its own MMCM for its 4\x clocks, so the host can move each channel's CK phase on its
own (\cref{sec:ddr}). The \textbf{accel} clock comes from a PLL on the second 200\,MHz oscillator and is set by the
build configuration (\code{clocks.accel\_mhz}); the kernel and its AXI side run on it. The PCIe hard block and the
PHY run on clocks derived from the PCIe reference. All domains are false-pathed against each other, and data
crosses only through the places marked in the figure: the PHY's asynchronous FIFOs, the asynchronous FIFOs inside
the LiteDRAM crossbar ports for the kernel, and synchronisers for the kernel's control registers. A frequency meter
reports the measured accel clock to the host, which uses it to convert cycles to seconds.

\begin{figure}[t]
  \centering
  \resizebox{\linewidth}{!}{% Clock sources, generated clocks, domains and the clock-domain crossings.
\begin{tikzpicture}[x=1cm,y=1cm, every node/.append style={font=\sffamily\footnotesize}]
  \tikzset{src/.style={blk, fill=white, minimum width=24mm},
           gen/.style={blk, fill=black!7, minimum width=15mm},
           dom/.style={blk, minimum width=50mm, align=left, text width=48mm},
           cdc/.style={draw=cRed!80!black, fill=cRed!8, rounded corners=1pt, font=\sffamily\scriptsize,
                       align=center, inner sep=2pt}}
  % sources
  \node[src] (o0) at (0, 0)    {oscillator \#0\\[-1pt]\scriptsize 200\,MHz (LVDS)};
  \node[src] (o1) at (0,-2.6)  {oscillator \#1\\[-1pt]\scriptsize 200\,MHz (LVDS)};
  \node[src] (op) at (0,-4.7)  {PCIe reference\\[-1pt]\scriptsize 100\,MHz from the slot};
  \node[src] (o5) at (0,-6.8)  {oscillator\\[-1pt]\scriptsize 50\,MHz};
  % generators
  \node[gen] (mm)  at (3.1, 0)    {S7MMCM};
  \node[gen] (pll) at (3.1,-2.6)  {S7PLL};
  \node[gen] (gtx) at (2.9,-4.7)  {GTX\\[-1pt]\scriptsize TXOUTCLK};
  \node[gen] (pm)  at (4.6,-4.7)  {MMCM};
  \node[gen] (dm)  at (6.2, 1.25) {2\x MMCM\\[-1pt]\scriptsize (per channel)};
  \draw[link] (o0) -- (mm);  \draw[link] (o1) -- (pll);
  \draw[link] (op) -- (gtx); \draw[link] (gtx) -- (pm);
  % domains
  \node[dom, fill=dSys] (sys) at (9.6, 0)
     {\textbf{sys} 125\,MHz\\[-1pt]\scriptsize CSR bus, LitePCIe endpoint + DMA, DDR DMA bridge, LiteDRAM crossbars + controllers};
  \node[dom, fill=dDdr] (ddr) at (9.6, 1.6)
     {\textbf{DDR PHY} 125 + 500\,MHz\\[-1pt]\scriptsize 4\x serdes = 1000\,MT/s; CK phase stepped by the host};
  \node[dom, fill=dAccel] (acc) at (9.6,-2.6)
     {\textbf{accel} 150 / 200 / 225 / 250\,MHz\\[-1pt]\scriptsize HLS kernel, AXI\(\rightarrow\)native, perf counters; from \code{clocks.accel\_mhz}};
  \node[dom, fill=dPcie] (pci) at (9.6,-4.7)
     {\textbf{PCIe} 125 / 250 / 500\,MHz\\[-1pt]\scriptsize pipe 125, \code{userclk2} 250 (128\,b AXIS), \code{userclk1} 500 (hard-block core, Gen2 x8)};
  \node[dom, fill=black!4] (fm) at (9.6,-6.8)
     {\textbf{clk50}\\[-1pt]\scriptsize frequency meters: accel and clk50 measured against sys};
  \node[blk, fill=black!4, minimum width=18mm] (idl) at (6.2,-1.15) {IDELAYCTRL\\[-1pt]\scriptsize 200\,MHz ref.};
  \draw[link] (mm) -- (sys);
  \draw[link] (mm.north) |- (dm.west);
  \draw[link] (dm) -- (ddr.west |- dm);
  \draw[link] (mm.south) |- (idl.west);
  \draw[link] (pll) -- (acc);
  \draw[link] (pm) -- (pci);
  \draw[link] (o5) -- (fm);
  % crossings: one lane per crossing, to the right of the domains
  \node[cdc, text width=31mm] (x1) at (14.25,-1.3) {\textbf{CDC accel \(\leftrightarrow\) sys}\\ crossbar ports: async FIFOs;\\ CSRs: PulseSynchronizer\\ (start), MultiReg (args,\\ status, counters)};
  \node[cdc, text width=27mm] (x2) at (17.35,-3.1) {\textbf{CDC PCIe \(\leftrightarrow\) sys}\\ PHY RX/TX datapath:\\ async FIFOs, 128\,b};
  \draw[<->, thick, cRed!80!black] ([yshift=-1.5mm]sys.east) -| (x1.north);
  \draw[<->, thick, cRed!80!black] (x1.south) |- (acc.east);
  \draw[<->, thick, cRed!80!black] ([yshift=2mm]sys.east) -| (x2.north);
  \draw[<->, thick, cRed!80!black] (x2.south) |- (pci.east);
  \node[note, anchor=north, text width=50mm] at (9.6,-7.45)
     {All domains are false-pathed against each other in Vivado; only the crossings above move data between them.};
\end{tikzpicture}
}
  \caption{Clock sources, the clock domains and the clock-domain crossings (CDC).}
  \label{fig:clocks}
\end{figure}
